\documentclass[a4paper,11pt]{article}

% =============================================================================
%  Charte graphique des documents LaTeX du projet (gabarit commenté).
%  Le préambule fixe la mise en forme ; le corps montre chaque élément une fois
%  et se remplace par le contenu du document. Un document reprend le préambule
%  sans modification, hors des passages bornés par les deux lignes
%  « % >>> À ADAPTER : <objet> » et « % <<< À ADAPTER » (comparaison par
%  script : à créer, ADR 0003 annoté le 7 novembre 2026). Le bloc « Macros du dépôt »,
%  en fin de préambule, porte \src et les rubriques de la fiche modèle.
%  Version sans : babel-french, caption, csquotes, dsfont, enumitem, float,
%  mathrsfs, mdframed, newtxtext, pdflscape, tikz/pgf, titlesec, tocloft.
%  Moteur : pdfLaTeX. Compatible avec les distributions anciennes
%  (MiKTeX 2.9, TeX Live 2012 et suivantes) : aucun crochet du noyau récent.
% =============================================================================

\usepackage[utf8]{inputenc}
\usepackage[T1]{fontenc}
% \headheight est laissé à 12 pt : fancyhdr le signale et l'agrandit lui-même
% (avertissement « \headheight is too small »), exactement comme dans le
% document source ; le fixer ici décalerait l'en-tête ou le texte de 1,6 pt.
\usepackage[left=2.5cm, right=2.5cm, top=2.5cm, bottom=2.5cm]{geometry}
\usepackage{setspace}
\usepackage{graphicx}
\usepackage{etoolbox}
\usepackage{keyval}
\usepackage{fancyhdr}
\usepackage{amsmath,amssymb,amsfonts}
\usepackage{url}
\usepackage[table]{xcolor}  % [table] : \rowcolors sur les anciennes versions
\usepackage{array}
\usepackage{booktabs}
\usepackage{colortbl}
\usepackage{longtable}
\usepackage{lscape}   % remplace pdflscape (rotation à l'écran : voir plus bas)
\usepackage{framed}   % remplace mdframed (encadré sécable)
\usepackage{atbegshi}  % crochets de sortie de page (remplace \AddToHook)
\usepackage[unicode]{hyperref}

% ------------------------------------------------------------- polices -------
% newtxtext composait le texte en TeX Gyre Termes X et le code en txtt ; les
% mathématiques restaient en Computer Modern. On retrouve la même chose sans
% newtxtext : TeX Gyre Termes (tgtermes) pour le texte, repli sur le Times de
% PSNFSS (ptm) s'il est absent ; txtt (txfonts) pour \ttfamily.
\IfFileExists{tgtermes.sty}
  {\usepackage{tgtermes}}
  {\renewcommand{\rmdefault}{ptm}}
\IfFileExists{t1txtt.fd}{\renewcommand{\ttdefault}{txtt}}{}

% ---------------------------------------------------------------- charte ----
% Bleu financier pour la structure, vert pour l'accent et les sorties, deux
% gris de tableau.
\definecolor{maincolor}{RGB}{0,70,140}
\definecolor{greenline}{RGB}{0,180,80}
\definecolor{tableheader}{RGB}{200,220,250}
\definecolor{tablerow}{RGB}{245,245,245}
\definecolor{financeblue}{RGB}{0,70,140}

% >>> À ADAPTER : métadonnées du document (page de garde, pied de page, PDF)
% Dans un guide, \datedoc est une date figée (jamais \today) et \statutdoc
% vaut Projet ou Final ; \datedroit est la date de référence du droit.
\newcommand{\annee}{2026}
\newcommand{\datedoc}{\today}
\newcommand{\versiondoc}{1.0}
\newcommand{\statutdoc}{Projet}
\newcommand{\titredoc}{Titre court du document}
\newcommand{\auteurdoc}{Prénom Nom}
\newcommand{\datedroit}{1\textsuperscript{er} octobre 2026}

\hypersetup{colorlinks=true, linkcolor=financeblue, urlcolor=financeblue,
            citecolor=financeblue, pdftitle={Titre du document},
            pdfauthor={\auteurdoc}}
% <<< À ADAPTER

\makeatletter

% ================================================ français sans babel ========
% Reprise, à la main, des effets de babel-french visibles dans le document
% source.
% (1) Césure française (motifs chargés dans le format pdfLaTeX).
\ifdefined\l@french
  \language=\l@french
\else
  \PackageWarningNoLine{charte}{Motifs de césure français absents :
    césure anglaise utilisée}
\fi
\lefthyphenmin=2 \righthyphenmin=2
% (2) Espacement « à la française » après les points.
\frenchspacing
% (3) Libellés.
\renewcommand{\contentsname}{Table des matières}
\renewcommand{\listfigurename}{Table des figures}
\renewcommand{\listtablename}{Liste des tableaux}
\renewcommand{\refname}{Références}
\renewcommand{\indexname}{Index}
\renewcommand{\figurename}{Figure}
\renewcommand{\tablename}{Table}
\renewcommand{\partname}{Partie}
\renewcommand{\appendixname}{Annexe}
\renewcommand{\abstractname}{Résumé}
% (4) Date en français (« 1er janvier 2026 »).
\renewcommand{\today}{\number\day\ifnum\day=1 \textsuperscript{er}\fi\space
  \ifcase\month\or janvier\or février\or mars\or avril\or mai\or juin\or
  juillet\or août\or septembre\or octobre\or novembre\or décembre\fi
  \space\number\year}
% (5) Premier paragraphe indenté après un titre.
\let\@afterindentfalse\@afterindenttrue
\@afterindenttrue
% (6) Listes : espacement vertical resserré de babel-french, appliqué à tous
%     les environnements construits sur \list.
\let\charte@listORI\list
\newdimen\charte@parskip
\def\charte@listV{%
  \setlength{\topsep}{0.8ex plus 0.4ex minus 0.4ex}%
  \setlength{\partopsep}{0.4ex plus 0.2ex minus 0.2ex}%
  \charte@parskip=\parskip
  \addtolength{\topsep}{-\charte@parskip}%
  \addtolength{\partopsep}{\charte@parskip}%
  \setlength{\itemsep}{0.4ex plus 0.2ex minus 0.2ex}%
  \setlength{\parsep}{0.4ex plus 0.2ex minus 0.2ex}%
  \ifdim\charte@parskip>\z@
    \setlength{\parsep}{\charte@parskip}%
    \addtolength{\itemsep}{-\charte@parskip}%
  \fi}
\def\list#1#2{\charte@listORI{#1}{\charte@listV #2}}
% (7) Ponctuation haute : sans babel-french, l'espace qui précède « : » et
%     l'espace fine qui précède « ; ! ? » ne sont plus ajoutées. On saisit
%     « ~: » (espace insécable) et « \fsp; », « \fsp! », « \fsp? » (espace
%     fine insécable, de même largeur que celle de babel-french).
\newcommand{\fsp}{\unskip\nobreak\hskip .5\fontdimen2\font\relax}
% (8) Puces en tiret cadratin à tous les niveaux.
\renewcommand{\labelitemi}{\textemdash}
\renewcommand{\labelitemii}{\textemdash}
\renewcommand{\labelitemiii}{\textemdash}
\renewcommand{\labelitemiv}{\textemdash}

% ================================================ listes sans enumitem =======
% itemize et enumerate acceptent un argument optionnel à clés, dont la
% syntaxe est celle d'enumitem pour les clés usuelles :
%   \begin{itemize}[itemsep=3pt] ... \end{itemize}
% Clés reconnues : itemsep, topsep, parsep, partopsep, leftmargin,
% labelsep, listparindent.
\define@key{charteliste}{itemsep}{\setlength{\itemsep}{#1}}
\define@key{charteliste}{topsep}{\setlength{\topsep}{#1}}
\define@key{charteliste}{parsep}{\setlength{\parsep}{#1}}
\define@key{charteliste}{partopsep}{\setlength{\partopsep}{#1}}
\define@key{charteliste}{leftmargin}{\setlength{\leftmargin}{#1}}
\define@key{charteliste}{labelsep}{\setlength{\labelsep}{#1}}
\define@key{charteliste}{listparindent}{\setlength{\listparindent}{#1}}

\renewenvironment{itemize}[1][]{%
  \ifnum\@itemdepth>\thr@@\@toodeep\else
    \advance\@itemdepth\@ne
    \edef\@itemitem{labelitem\romannumeral\the\@itemdepth}%
    \expandafter\list\csname\@itemitem\endcsname
      {\def\makelabel##1{\hss\llap{##1}}\setkeys{charteliste}{#1}}%
  \fi}{\endlist}

\renewenvironment{enumerate}[1][]{%
  \ifnum\@enumdepth>\thr@@\@toodeep\else
    \advance\@enumdepth\@ne
    \edef\@enumctr{enum\romannumeral\the\@enumdepth}%
    \expandafter\list\csname label\@enumctr\endcsname
      {\usecounter\@enumctr\def\makelabel##1{\hss\llap{##1}}%
       \setkeys{charteliste}{#1}}%
  \fi}{\endlist}

% ================================================ légendes sans caption ======
% Étiquette en gras, couleur maincolor, petites capitales ; séparateur
% « : » précédé d'une espace insécable, comme dans le document source.
\renewcommand{\fnum@figure}{{\scshape\figurename}\nobreakspace\thefigure}
\renewcommand{\fnum@table}{{\scshape\tablename}\nobreakspace\thetable}
\newcommand{\charte@etiquette}[1]{%
  {\bfseries\color{maincolor}#1\nobreakspace:}\space}
% Comme avec caption : légende composée en interligne simple, entre deux
% étais, sans ressort interligne au-dessus ; centrée si elle tient sur une
% ligne, justifiée sinon.
\long\def\@makecaption#1#2{%
  \vskip\abovecaptionskip
  \sbox\@tempboxa{\normalfont\charte@etiquette{#1}#2}%
  \nointerlineskip
  \ifdim\wd\@tempboxa>\hsize
    \hb@xt@\hsize{\vbox{\setstretch{1}\normalfont
      \parindent\z@ \leftskip\z@ \rightskip\z@ \parfillskip\z@\@plus1fil
      \strut\charte@etiquette{#1}#2\strut\par}}%
  \else
    \global\@minipagefalse
    \hb@xt@\hsize{\hfil\setstretch{1}\normalfont\strut\box\@tempboxa\hfil}%
  \fi
  \vskip\belowcaptionskip}

% Figures et tableaux « ici » : reprise exacte de la position [H] du
% package float (boîte non flottante, \intextsep avant et après).
\newsavebox\charte@boiteH
\def\charte@debutH#1{%
  \def\@captype{#1}%
  \global\setbox\charte@boiteH\color@vbox\normalcolor
    \vbox\bgroup\hsize\columnwidth\@parboxrestore
      \@floatboxreset\@setnobreak\ignorespaces}
\def\charte@finH{%
  \@endfloatbox
  \vskip\intextsep\box\charte@boiteH\vskip\intextsep\relax}
\newenvironment{figureH}{\charte@debutH{figure}}{\charte@finH}
\newenvironment{tableH}{\charte@debutH{table}}{\charte@finH}

% ================================================ en-têtes et pieds ==========
\renewcommand{\baselinestretch}{1.08}
\pagestyle{fancy}
\fancyhf{}
\fancyhead[L]{\color{maincolor}\nouppercase{\leftmark}}
\renewcommand{\headrulewidth}{0pt}
\renewcommand{\footrulewidth}{0.4pt}
\fancyfoot[R]{\color{greenline}\textbf{page \thepage}}
\fancyfoot[L]{\color{maincolor}\titredoc}
\fancypagestyle{plain}{\fancyhf{}%
  \renewcommand{\headrulewidth}{0pt}\renewcommand{\footrulewidth}{0.4pt}%
  \fancyfoot[R]{\color{greenline}\textbf{page \thepage}}%
  \fancyfoot[L]{\color{maincolor}\titredoc}}

% ================================================ titres sans titlesec =======
% Bandeau : filet bleu pleine largeur, titre vert en dessous. Utilisé seul
% pour les pages non numérotées (objet, sommaire) et par \section.
\newcommand{\titrebandeau}[1]{%
  \noindent{\color{financeblue}\rule{\linewidth}{2.5pt}}
  \par\vspace{5mm}
  {\color{greenline}\raggedright\fontsize{19}{23}\bfseries\selectfont #1\par}
  \vspace{6mm}}

% Numéros de titre : aucun pour \section (le numéro ne figure qu'au
% sommaire et dans l'en-tête), 0,7 em puis 0,6 em d'écart ensuite.
\def\@seccntformat#1{\csname charte@num@#1\endcsname}
\def\charte@num@section{}
\def\charte@num@subsection{\strut\thesubsection\hspace{0.7em}}
\def\charte@num@subsubsection{\strut\thesubsubsection\hspace{0.6em}}
\def\charte@num@paragraph{\theparagraph\quad}
\def\charte@num@subparagraph{\thesubparagraph\quad}

% \@startsection passe le titre, entre accolades, à la dernière commande de
% son argument de style. Pour \section, cette commande est \charte@bandeau :
% une ligne vide (étai) au-dessus du bandeau et une autre en dessous, et un
% étai dans chaque ligne, comme la forme « display » de titlesec avec
% étiquette vide.
\newcommand{\charte@bandeau}[1]{%
  \parindent\z@ \everypar{\strut}%
  {\strut\par}\nobreak
  \titrebandeau{#1}\nobreak
  \kern\z@\nobreak\strut\par}
% Titres en ligne : un étai en fin de titre (le premier est dans le numéro).
\newcommand{\charte@etai}[1]{\charte@etai@#1}
\def\charte@etai@#1\@@par{#1\kern\z@\strut\@@par}
% Réserve avant un titre : si le bas de la page ne peut plus recevoir le titre
% et trois lignes de texte, le titre passe à la page suivante (M41 : aucun
% titre en bas de page avec moins de trois lignes sous lui). La réserve
% compte le titre lui-même et ses espaces (mesurés : bandeau de section
% 7,4 lignes, sous-section 2,9 lignes), trois lignes et le cadre d'un
% encadré qui suivrait : 12 lignes pour \section, 7 pour les autres. Logique de la
% macro \Needspace* du package needspace, sans le charger : la place restante
% est mesurée (\pagegoal - \pagetotal, après une pénalité nulle qui met la
% page à jour) et le saut n'est forcé que si elle manque. La forme élastique
% de \needspace, employée jusqu'ici, créait un point de coupure de faible
% coût que TeX préférait à toute autre coupure dès qu'un bloc insécable
% suivait (tableau, encadré), même avec la page aux trois quarts vide. Sans
% effet juste après un autre titre (\if@nobreak), où elle créerait un point
% de coupure sous ce titre, ni dans une boîte (\ifinner : minipage, \parbox),
% où la page en cours ne la concerne pas.
\newcommand{\charte@reserve}[1]{\par\ifinner\else\if@nobreak\else
  \penalty\z@ % met à jour \pagetotal
  \begingroup\setlength{\dimen@}{#1}%
  \dimen@ii\pagegoal \advance\dimen@ii-\pagetotal
  \ifdim\pagegoal=\maxdimen\else\ifdim\dimen@>\dimen@ii
    \vfil\penalty-\@M % place insuffisante : page suivante
  \fi\fi\endgroup\fi\fi}
\renewcommand{\section}{\charte@reserve{12\baselineskip}%
  \@startsection{section}{1}{\z@}%
  {\z@}{1sp}{\normalfont\charte@bandeau}}
% Espacements : ceux de titlesec, soit 3,25 ex (+1 ex, -0,2 ex) avant et
% 1,5 ex (+0,2 ex) après, l'ex étant celui que titlesec mémorise au
% chargement (Computer Modern à 11 pt, soit 4,71341 pt).
\newdimen\charte@ex \charte@ex=4.71341pt
\newskip\charte@avanttitre
\charte@avanttitre=-3.25\charte@ex \@plus -1\charte@ex \@minus -.2\charte@ex
\newskip\charte@aprestitre
\charte@aprestitre=1.5\charte@ex \@plus .2\charte@ex
\renewcommand{\subsection}{\charte@reserve{7\baselineskip}%
  \@startsection{subsection}{2}{\z@}%
  {\charte@avanttitre}{\charte@aprestitre}%
  {\normalfont\large\bfseries\color{financeblue}\charte@etai}}
\renewcommand{\subsubsection}{\charte@reserve{7\baselineskip}%
  \@startsection{subsubsection}{3}{\z@}%
  {\charte@avanttitre}{\charte@aprestitre}%
  {\normalfont\normalsize\bfseries\color{financeblue}\charte@etai}}

% ================================================ sommaire sans tocloft ======
% Sections : gras bleu, numéro sur 2,2 em, numéro de page bleu (maigre),
% sans points de conduite. Sous-sections : retrait 1,5 em, numéro sur 2,3 em,
% points de conduite. L'em est celui que tocloft mémorisait au chargement
% (10,88788 pt), pour un sommaire identique au document source.
\newdimen\charte@em \charte@em=10.88788pt
\renewcommand*{\l@section}[2]{%
  \ifnum\c@tocdepth>\z@
    \addpenalty\@secpenalty
    \addvspace{1.0\charte@em \@plus\p@}%
    \setlength\@tempdima{2.2\charte@em}%
    \begingroup
      \parindent\z@ \rightskip\@pnumwidth
      \parfillskip -\@pnumwidth
      \leavevmode
      \advance\leftskip\@tempdima
      \hskip -\leftskip
      {\bfseries\color{financeblue}#1}\nobreak\hfil\nobreak
      \hb@xt@\@pnumwidth{\hss\color{financeblue}#2}\par
    \endgroup
  \fi}
\renewcommand*{\l@subsection}{\@dottedtocline{2}{1.5\charte@em}{2.3\charte@em}}
\setcounter{tocdepth}{2}

% ================================================ environnements de texte ====
\newcommand{\code}[1]{\texttt{\small #1}}
\newcommand{\refl}[1]{{\small\textcolor{maincolor!70!black}{\textit{Référence.} #1}}\par\smallskip}

% Encadré de citation réglementaire : filet vertical bleu en marge.
\newcommand{\reglement}[2]{%
  \par\vspace{2mm}\noindent\begin{minipage}{\linewidth}
  \small\setlength{\parindent}{0pt}
  \textcolor{financeblue}{\vrule width 2pt}\hspace{4mm}%
  \begin{minipage}{\dimexpr\linewidth-7mm\relax}
  \textcolor{financeblue}{\textbf{#1}}
  \par\vspace{2pt}
  {\itshape #2}
  \end{minipage}
  \end{minipage}\par\vspace{3mm}}

% Encadré générique : fond gris, marge intérieure 6 mm à gauche et 10 pt à
% droite, 4 pt en haut et en bas, 3 mm avant (sauf après un titre) et
% après ; premier paragraphe non
% indenté. Sécable entre pages (framed), selon \charte@coupe ci-dessous. Le
% cadre garde la profondeur de la boîte d'essai de framed : sa hauteur
% mesurée est de 8 pt (elle valait -13,68 pt quand la profondeur était
% annulée, et framed taillait le premier morceau trop grand pour la page).
\newenvironment{encadre}{%
  \def\FrameCommand##1{{\setlength{\fboxsep}{\z@}%
    \colorbox{tablerow}{\hspace*{6mm}%
      \vbox{\vskip4pt\hbox{##1}\vskip4pt}\hspace*{10pt}}}}%
  \@ifundefined{OuterFrameSep}{}{\setlength{\OuterFrameSep}{\z@}}%
  \topsep\z@ \partopsep\z@ % framed antérieur à 0.96 : espacement via \topsep
  \par
  \if@nobreak\charte@sanscoupure\else\vskip 3mm\fi
  \nointerlineskip
  \MakeFramed{\let\penalty\charte@penalite@ori
    \advance\hsize-\width \FrameRestore
    \everypar{{\setbox\z@\lastbox}\everypar{}}}}%
  {\global\let\fb@put@frame\charte@coupe % hors du groupe de la boîte
   \global\charte@premiertrue
   \endMakeFramed
   \global\let\fb@put@frame\charte@putframe
   \par
   \vskip 3mm\prevdepth 4pt
   \global\@nobreakfalse\global\@endpetrue}

% Coupe de l'encadré (E3.7, M41), décidée à chaque placement d'un morceau,
% avant la coupe de framed. Premier placement : un encadré court (hauteur
% totale au plus \charte@encadrecourt) qui ne tient pas dans la place
% restante est posé entier, et TeX coupe la page au dernier point de coupure
% qui le précède : juste avant lui, ou avant le titre qu'il suit (aucun point
% de coupure entre un titre et l'encadré : \charte@sanscoupure). Encadré
% long : coupé, chaque morceau gardant au moins trois lignes, à chaque page
% (page suivante si la place restante en tient moins ; page écourtée si le
% reste en tiendrait moins). Exception : si des flottants attendent la page
% suivante (\@deferlist) et que l'encadré court ne tient pas avec eux sur une
% page, posé entier il déborderait ; il est alors coupé comme un encadré
% long (M200). De même s'il commence en tête de page : il ne peut aller plus
% loin, et la page peut déjà porter des flottants (M202). Un encadré coupé trop
% court pour donner deux morceaux de trois lignes (hauteur du contenu, cadre
% déduit, sous 5,5 interlignes : six lignes ou moins ; sept lignes se coupent
% en trois et quatre) passe entier à la page suivante. Limites : tous les
% flottants en attente sont comptés, même ceux qui n'iront pas sur la page
% suivante ([p], [b] trop grand, \topnumber dépassé) : coupe parfois
% inutile, avertie par le contrôle 13 ; sur une page qui porte un flottant
% en tête, la coupe de framed peut déborder d'environ 1 pt (Overfull \vbox
% mesuré jusqu'à 1,03 pt, toléré en AVERT par le contrôle 3 jusqu'à 1,5 pt,
% M202).
\newdimen\charte@encadrecourt \charte@encadrecourt=.5\textheight
\newif\ifcharte@premier
\let\charte@putframe\fb@put@frame
\let\charte@penalite@ori\penalty
% Pénalités 9999 de \MakeFramed (mise à jour de \pagetotal) rendues
% insécables juste après un titre ; \penalty est rétabli dans la boîte.
\def\charte@sanscoupure{\def\penalty{\afterassignment\charte@penalite\count@}}
\def\charte@penalite{\ifnum\count@=9999 \count@\@M\fi\charte@penalite@ori\count@}
\def\charte@coupe#1#2{%
  \let\charte@suite\charte@putframe
  \ifinner\else\ifdim\pagegoal=\maxdimen\else
    \fb@sizeofframe#1%
    \dimen@\pagegoal \advance\dimen@-\pagetotal \advance\dimen@-\fb@frh
    \advance\dimen@-\pageshrink \advance\dimen@-\pagedepth
    \dimen@ii\ht\@tempboxa
    \ifdim\dimen@ii>\dimen@
      \advance\dimen@ii\fb@frh
      \ifcharte@premier\ifdim\dimen@ii>\charte@encadrecourt\else
        \charte@flottants % hauteur des flottants en attente (M200)
        \advance\dimen@ii\@tempdima
        \@tempdimb\topskip \advance\@tempdimb\p@
        \ifdim\pagetotal>\@tempdimb % la page porte déjà du texte
          \ifdim\dimen@ii>\textheight\else % tient avec eux : posé entier
            \let\charte@suite\charte@entier
          \fi
        \fi % en tête de page (flottants déjà placés) : coupé (M202)
        \advance\dimen@ii-\@tempdima
      \fi\fi
      \advance\dimen@ii-\fb@frh
      \ifx\charte@suite\charte@entier\else
        \ifdim\dimen@<3\baselineskip
          \eject\fb@adjheight
        \else
          \@tempdimb\dimen@ii \advance\@tempdimb-\fb@frh
          \advance\dimen@ii-\dimen@
          \ifdim\dimen@ii<3\baselineskip
            \ifdim\@tempdimb<5.5\baselineskip
              % trop court pour deux morceaux de trois lignes : page suivante
              \eject\fb@adjheight
            \else
              \advance\dimen@ii-3\baselineskip
              \advance\dimen@ii-\fb@frh % marge : rembourrage du cadre
              \advance\dimen@ii\pagegoal \pagegoal\dimen@ii % toujours global
            \fi
          \fi
        \fi
      \fi
    \fi
  \fi\fi
  \global\charte@premierfalse
  \charte@suite#1#2}
% Hauteur (\@tempdima) des flottants qui attendent la page suivante
% (\@deferlist), avec leurs espaces ; zéro s'il n'y en a pas.
\def\charte@flottants{\@tempdima\z@
  \ifx\@deferlist\@empty\else
    \begingroup
      \def\@elt##1{\global\advance\@tempdima\ht##1%
        \global\advance\@tempdima\dp##1\global\advance\@tempdima\floatsep}%
      \@deferlist
    \endgroup
    \advance\@tempdima\textfloatsep \advance\@tempdima-\floatsep
  \fi}
% Encadré court posé entier : la page déborde et TeX la coupe avant lui.
\def\charte@entier#1#2{\fb@putboxa#1\fb@afterframe}

% Fiches : un environnement, deux jeux de rubriques fixes. Fiche méthode
% (corps du guide) : six rubriques, ouvertes par \Principe, ci-dessous ;
% fiche modèle (annexe technique) : cinq rubriques, ouvertes par \Cadre, au
% bloc « Macros du dépôt ».
\newcounter{ficheid}
\newenvironment{fiche}[1]{%
  \stepcounter{ficheid}%
  \phantomsection
  \subsection{#1}%
  \label{fiche:\theficheid}%
}{%
  \end{rubriques}\vspace{0.5em}%
}
% Liste des rubriques : aucun retrait ; l'intitulé fait partie du
% paragraphe (ses espaces se justifient avec la ligne, comme le style
% « unboxed » d'enumitem), suivi d'un écart fixe \labelsep.
\newenvironment{rubriques}{%
  \list{}{\setlength{\leftmargin}{\z@}\setlength{\labelwidth}{\z@}%
          \setlength{\itemindent}{\z@}%
          \setlength{\itemsep}{3pt}\setlength{\topsep}{3pt}%
          \setlength{\parsep}{2pt}%
          \let\makelabel\@gobble}%
  \clubpenalties 3 10000 10000 0
  \@beginparpenalty\@M \global\charte@rubriques\z@}{\endlist}
% Un titre de rubrique n'est jamais seul en bas de page : il fait partie de
% son premier paragraphe, dont les deux premières coupures sont interdites
% (\clubpenalties, liste des rubriques), et une liste qui le suit ne s'en
% sépare pas (\@beginparpenalty) ; aucune réserve de place, qui laisserait
% orphelin le titre de la fiche ; aucune coupure avant la deuxième ni la
% troisième rubrique, pour que le titre de la fiche garde trois lignes sous
% lui même quand les premières rubriques tiennent sur une ligne. Le réglage
% \clubpenalties vaut pour tous les paragraphes de la fiche, listes
% imbriquées comprises : aucun ne laisse moins de trois lignes en bas de page
% (un paragraphe court passe alors entier à la page suivante).
% Limite : un titre de rubrique suivi d'une liste peut garder une seule
% puce sous lui en bas de page (jamais aucune) ; la règle des trois lignes
% (M41) est celle des titres de section (M202).
\newcount\charte@rubriques
\newcommand{\rubrique}[1]{%
  \global\advance\charte@rubriques\@ne
  \ifnum\charte@rubriques<4 \@itempenalty\@M\else\@itempenalty-\@lowpenalty\fi
  \item[]\@itempenalty-\@lowpenalty % listes de la rubrique : sécables
  \textcolor{financeblue}{\textbf{#1}}\hskip\labelsep\ignorespaces}
\newcommand{\Principe}{%
  \begin{rubriques}%
  \rubrique{1. Principe et modèle de référence.}}
\newcommand{\Donnees}{\rubrique{2. Données requises.}}
\newcommand{\Conditions}{\rubrique{3. Conditions d'usage.}}
\newcommand{\Forces}{\rubrique{4. Forces.}}
\newcommand{\Limites}{\rubrique{5. Limites et pièges.}}
\newcommand{\Referentiels}{\rubrique{6. Référentiels concernés.}}

% Titre de partie : un filet bleu, « Partie » en petites capitales, titre
% vert.
\newcommand{\partie}[2]{%
  \clearpage
  \vspace*{1cm}
  \noindent{\color{financeblue}\rule{\linewidth}{2pt}}\par\vspace{3mm}
  {\color{financeblue}\raggedright\normalsize\bfseries\scshape #1\par}
  \vspace{1.5mm}
  {\color{greenline}\raggedright\fontsize{20}{24}\bfseries\selectfont #2\par}
  \vspace{5mm}
  \addcontentsline{toc}{section}{\protect\numberline{}\textsc{#1}~: #2}}

% ================================================ page de garde sans TikZ ====
% \charte@surpage{<matériel vertical>} superpose, sur la prochaine page
% expédiée, un contenu placé depuis le coin supérieur gauche du papier
% (atbegshi : fonctionne avec les noyaux LaTeX anciens comme récents).
\newcommand{\charte@surpage}[1]{%
  \AtBeginShipoutNext{%
    \global\setbox\AtBeginShipoutBox\vbox{\offinterlineskip
      \vbox to\z@{\kern-1in\kern-\voffset
        \hbox to\z@{\kern-1in\kern-\hoffset
          \vbox to\z@{#1\vss}\hss}\vss}%
      \box\AtBeginShipoutBox}}}
% Trois filets : bleu de 3,2 à 3,28 cm du haut avec amorce verte de 5 cm,
% bleu de 3,6 à 3,53 cm du bas, entre les marges de 2,5 cm.
% \leavevmode ouvre le paragraphe comme le faisait la figure TikZ du document
% source : la mention d'en-tête se place ainsi au-dessus du filet, en retrait.
\newcommand{\filetsgarde}{%
  \charte@surpage{%
    \kern3.2cm
    \moveright2.5cm\hbox{\color{financeblue}%
      \rule{\dimexpr\paperwidth-5cm\relax}{0.08cm}}%
    \kern-0.08cm
    \moveright2.5cm\hbox{\color{greenline}\rule{5cm}{0.08cm}}%
    \kern\dimexpr\paperheight-3.6cm-3.28cm\relax
    \moveright2.5cm\hbox{\color{financeblue}%
      \rule{\dimexpr\paperwidth-5cm\relax}{0.07cm}}}%
  \leavevmode}

% ================================================ pages en paysage ===========
% lscape tourne le contenu ; comme pdflscape, on ajoute /Rotate 90 aux pages
% concernées pour qu'elles s'affichent à l'endroit à l'écran.
\newif\ifcharte@paysage
\apptocmd{\landscape}{\global\charte@paysagetrue}{}{}
\apptocmd{\endlandscape}{\global\charte@paysagefalse}{}{}
\AtBeginShipout{%
  \ifcharte@paysage\global\pdfpageattr{/Rotate 90}%
  \else\global\pdfpageattr{}\fi}

% ================================================ Macros du dépôt ============
% Propres au dépôt de guides, identiques dans tous les documents.
%
% (1) Statuts de vérification des sources. statuts_sources.tex est généré
%     depuis le registre des sources (sources/registre.csv) par
%     outils/controles.sh, et trouvé par TEXINPUTS (répertoire sources/) ;
%     il contient une ligne \charte@statut{ID}{statut} par identifiant, le
%     statut valant verifie, incertain, nontrouve ou absent.
\newif\ifcharte@statuts
\newcommand{\charte@statut}[2]{\@namedef{charte@src@#1}{#2}}
\def\charte@verifie{verifie}
\IfFileExists{statuts_sources.tex}
  {\input{statuts_sources.tex}\charte@statutstrue}
  {\PackageWarningNoLine{charte}{statuts_sources.tex introuvable :
     statuts de vérification des sources non calculés}}
% (2) \src{ID} : identifiant du registre des sources, placé dans le premier
%     argument de \reglement ou dans l'argument de \refl, en fin de texte.
%     Statut vérifié : rien d'imprimé ; incertain, non trouvé ou absent : une
%     marque † en exposant, accolée au mot qui précède ; identifiant inconnu :
%     [ID ?] en gras et un avertissement « Source inconnue au registre ».
\DeclareRobustCommand{\src}[1]{%
  \ifcharte@statuts
    \@ifundefined{charte@src@#1}%
      {\textbf{[#1~?]}%
       \PackageWarning{charte}{Source inconnue au registre : #1}}%
      {\unskip
       \expandafter\ifx\csname charte@src@#1\endcsname\charte@verifie
       \else\textsuperscript{\dag}\fi}%
  \fi
  \protected@write\@auxout{}{\string\charte@srcpage{#1}{\thepage}}}
% (2 bis) Folios des \src, pour l'annexe des sources citées.
\newcommand{\charte@srcpage}[2]{%
  \@ifundefined{charte@pg@#1}%
    {\global\@namedef{charte@pg@#1}{\hyperpage{#2}}%
     \global\@namedef{charte@dp@#1}{#2}}%
    {\def\charte@tmp{#2}%
     \expandafter\ifx\csname charte@dp@#1\endcsname\charte@tmp\else
       \expandafter\g@addto@macro\csname charte@pg@#1\endcsname{, \hyperpage{#2}}%
       \global\@namedef{charte@dp@#1}{#2}\fi}}
\AtEndDocument{\let\charte@srcpage\@gobbletwo}
\newcommand{\srcfolios}[1]{%
  \@ifundefined{charte@pg@#1}%
    {\textbf{??}\G@refundefinedtrue
     \PackageWarning{charte}{Folios de la source #1 non definis}}%
    {\@nameuse{charte@pg@#1}}}
% (3) Rubriques de la fiche modèle (annexe technique), sur l'environnement
%     fiche : cinq rubriques fixes, ouvertes par \Cadre.
\newcommand{\Cadre}{%
  \begin{rubriques}%
  \rubrique{1. Cadre et notations.}}
\newcommand{\Hypotheses}{\rubrique{2. Hypothèses.}}
\newcommand{\Utilisation}{\rubrique{3. Conditions d'utilisation.}}
\newcommand{\Formules}{\rubrique{4. Formules et statut.}}
\newcommand{\LimitesModele}{\rubrique{5. Limites.}}
% (4) Bibliographie (BibTeX, style plain, base references, soit
%     sources/references.bib trouvé par BIBINPUTS) : \bibliographiedepot, en
%     fin de document, ne la compose que si le document cite une entrée par
%     \cite ou \nocite ; sans citation, bibtex s'arrêterait en erreur
%     (« I found no \citation commands »).
\newif\ifcharte@citations
\pretocmd{\@citex}{\global\charte@citationstrue}{}
  {\PackageWarningNoLine{charte}{\string\@citex\space non modifiée :
     bibliographie non composée}}
\pretocmd{\nocite}{\global\charte@citationstrue}{}
  {\PackageWarningNoLine{charte}{\string\nocite\space non modifiée :
     bibliographie non composée}}
\newcommand{\bibliographiedepot}{%
  \ifcharte@citations
    \bibliographystyle{plain}%
    \bibliography{references}%
  \fi}
% (5) Annexe des sources citées. \sourcescitees charge sources_citees.tex,
%     tableau long généré à chaque compilation par
%     outils/controles/sources_citees.R depuis le registre des sources et les
%     \src du document, et trouvé par TEXINPUTS (répertoire de construction) ;
%     l'annexe ne se saisit jamais à la main.
\newcommand{\sourcescitees}{%
  \IfFileExists{sources_citees.tex}
    {\input{sources_citees.tex}}
    {\PackageWarningNoLine{charte}{sources_citees.tex introuvable :
       annexe des sources citées non générée}}}

\makeatother

\emergencystretch=1.5em

\begin{document}

% ============================================================ page de garde ==
\thispagestyle{empty}
\filetsgarde
\vspace*{2.1cm}
\noindent
{\color{financeblue}\normalsize\bfseries\scshape Édition \annee}

\vspace{1.7cm}
\noindent
{\color{financeblue}\fontsize{28}{34}\bfseries\selectfont
 Titre principal\\ du document\par}

\vspace{1.1cm}
\noindent
{\color{black}\Large
 Sous-titre sur une première ligne\\
 Sous-titre sur une seconde ligne\par}

\vspace{0.7cm}
\noindent
{\color{black}\normalsize\itshape
 Mention de cadrage~: garantie, organismes d'assurance visés, référentiels\par}

\vfill
\noindent
{\footnotesize\color{black}
 Date de référence du droit~: \datedroit\par
 \smallskip\noindent
 Établie le \datedoc\quad\textbar\quad Version \versiondoc\quad\textbar\quad
 Statut~: \statutdoc\quad\textbar\quad Auteur~: \auteurdoc\par}
\vspace{2.6cm}
\clearpage

% ================================================================= objet ====
\titrebandeau{Objet du document}
\markboth{Objet du document}{}

\noindent
Ce gabarit reproduit la \textbf{charte graphique} de la documentation
du projet. Chaque élément de la charte y figure
une fois, dans son contexte d'usage~: page de garde, page d'objet, sommaire,
titres de partie, de section et de sous-section, fiche méthode et fiche
modèle, encadrés, citations réglementaires, tableaux, tableau long en paysage
et figure légendée.

\medskip\noindent
Il est organisé en trois parties. La \textbf{partie~I} présente la palette,
la typographie et les niveaux de titre. La \textbf{partie~II} montre les
blocs de texte (encadrés, citations, références, listes) et les tableaux. La
\textbf{partie~III} donne le modèle de fiche méthode, de fiche modèle et de
figure, puis la table de correspondance des packages remplacés.

\vspace{8mm}
\reglement{Code des assurances, article R343-7 \src{CA-R343-7}}{Emplacement
de l'extrait cité, en italique, précédé de son intitulé en gras bleu et bordé
d'un filet vertical en marge\fsp; l'intitulé se termine par l'identifiant du
registre des sources.}

\vspace{6mm}
\noindent
\begin{center}
\small
\rowcolors{2}{tablerow}{white}
\begin{tabular}{@{}p{5.2cm}p{9.6cm}@{}}
\rowcolor{tableheader}
\textcolor{financeblue}{\textbf{Élément}} & \textcolor{financeblue}{\textbf{Référence}} \\
Texte de référence & Intitulé du texte, articles concernés\fsp; annexes \\
Cahier des charges & \code{docs/exigences.md} \\
Module de calcul & \code{<module de calcul>} \\
Paramètre principal & $n = 100$ observations \\
\end{tabular}
\end{center}

\clearpage
\titrebandeau{Sommaire}
\markboth{Sommaire}{}
% \tableofcontents composerait son propre titre, en doublon avec le bandeau :
% on n'appelle que \@starttoc, qui compose les entrées seules.
\makeatletter
\@starttoc{toc}
\makeatother
\newpage


\partie{Partie I}{Identité visuelle}

% =============================================================================
\section{Palette et typographie}
\label{sec:palette}
% =============================================================================

\subsection{Couleurs}

La charte repose sur deux couleurs d'identité et deux gris de tableau. Le bleu
porte la structure (filets, titres de sous-section, liens, étiquettes), le
vert porte l'accent (titres de section et de partie, numéro de page, sorties).

\begin{center}
\small
\rowcolors{2}{tablerow}{white}
\begin{tabular}{@{}p{3.0cm}p{2.6cm}p{1.4cm}p{7.0cm}@{}}
\rowcolor{tableheader}
\textcolor{financeblue}{\textbf{Nom}} &
\textcolor{financeblue}{\textbf{RVB}} &
\textcolor{financeblue}{\textbf{Aperçu}} &
\textcolor{financeblue}{\textbf{Usage}} \\
\code{financeblue} & 0, 70, 140 & \textcolor{financeblue}{\rule{1cm}{2.2ex}} &
Filets, titres de sous-section, sommaire, liens, rubriques de fiche \\
\code{maincolor} & 0, 70, 140 & \textcolor{maincolor}{\rule{1cm}{2.2ex}} &
En-tête, pied de page, étiquettes de légende, références \\
\code{greenline} & 0, 180, 80 & \textcolor{greenline}{\rule{1cm}{2.2ex}} &
Titres de section et de partie, numéro de page, amorce de la page de garde \\
\code{tableheader} & 200, 220, 250 & \textcolor{tableheader}{\rule{1cm}{2.2ex}} &
Ligne d'en-tête des tableaux \\
\code{tablerow} & 245, 245, 245 & \textcolor{tablerow}{\rule{1cm}{2.2ex}} &
Lignes alternées des tableaux, fond des encadrés \\
\end{tabular}
\end{center}

\subsection{Typographie}

Le texte courant est composé en Termes (clone de Times), corps 11~pt,
interligne multiplié par 1,08\fsp; le code, en \code{txtt}, via la macro
\code{\textbackslash code\{\}}\fsp; les mathématiques restent en Computer Modern,
par exemple $\bar x = \frac{1}{n}\sum_{i=1}^{n} x_i$.

\subsubsection{Niveau de sous-sous-section}

Troisième niveau de titre, en gras bleu au corps du texte. Il n'apparaît pas
au sommaire (profondeur limitée à deux niveaux).

\subsection{Niveaux de titre}

\begin{itemize}[itemsep=2pt]
\item \textbf{Partie}~: \code{\textbackslash partie\{Partie I\}\{Titre\}},
  nouvelle page, filet bleu, mention en petites capitales, titre vert.
\item \textbf{Section}~: \code{\textbackslash section\{Titre\}}, bandeau
  (filet bleu épais, titre vert en 19~pt), numéro visible seulement au
  sommaire et dans l'en-tête.
\item \textbf{Sous-section et sous-sous-section}~: gras bleu, numérotées.
\item \textbf{Bandeau seul}~: \code{\textbackslash titrebandeau\{Titre\}}
  suivi de \code{\textbackslash markboth\{Titre\}\{\}} pour une page non
  numérotée (objet, sommaire).
\end{itemize}


\partie{Partie II}{Blocs de texte et tableaux}

% =============================================================================
\section{Encadrés, citations et listes}
\label{sec:blocs}
% =============================================================================

\subsection{Encadré générique}

L'encadré sert aux remarques de portée et aux extraits de code. Un encadré
court (au plus la moitié de la hauteur de page) ne se coupe pas~: s'il
ne tient pas en bas de page, il passe entier à la page suivante, sauf si
des flottants l'attendent ou s'il commence en tête de page sous un flottant
(il se coupe alors comme un encadré long, s'il est assez long pour deux
morceaux de trois lignes~; sinon il passe entier à la page suivante). Un encadré
plus long se coupe entre deux pages, chaque morceau gardant au moins trois
lignes.

\begin{encadre}
\small
\code{donnees <- lire\_donnees("chemin/des/donnees")}\\
\code{res <- calculer(donnees, methode = "standard", n = 100,}\\
\code{\phantom{res <- calculer(}graine = 20260831)}\\
\code{res\$resultat\_final \ \ ;\ \ tableau\_resultats(res)}
\end{encadre}

\begin{encadre}
\textbf{Portée.} L'encadré de portée dit ce que l'outil fait et ce qu'il ne
fait pas~: les cas du texte de référence qu'il ne couvre pas, et les méthodes
qu'il met en \oe uvre.
\end{encadre}

\subsection{Citation réglementaire et référence}

Toute citation, par \code{\textbackslash reglement} ou
\code{\textbackslash refl}, se termine par l'identifiant du registre des
sources, \code{\textbackslash src\{ID\}}. Rien ne s'imprime pour un fait
vérifié\fsp; une marque en exposant signale un fait incertain ou non trouvé
(référence ci-dessous)\fsp; un identifiant absent du registre s'imprime en gras
entre crochets.

\refl{Code des assurances, articles R351-x \src{CA-R351}}

\reglement{Code des assurances, article L351-2 \src{CA-L351-2}}{Emplacement
d'un second extrait.}

\subsection{Listes}

Liste à puces (tiret cadratin), avec un écart réglable par
\code{[itemsep=...]}~:

\begin{itemize}[itemsep=3pt]
\item premier élément, sur une ligne\fsp;
\item second élément, plus long, pour montrer l'alignement des lignes suivantes
  sous le texte de l'élément et non sous la puce\fsp;
  \begin{itemize}
  \item élément de second niveau.
  \end{itemize}
\item dernier élément.
\end{itemize}

Liste numérotée~:

\begin{enumerate}[itemsep=2pt]
\item contrôle des données d'entrée\fsp;
\item calcul\fsp;
\item restitution des résultats.
\end{enumerate}

% =============================================================================
\section{Tableaux}
\label{sec:tableaux}
% =============================================================================

\subsection{Tableau courant}

En-tête sur fond \code{tableheader} en gras bleu, lignes alternées
\code{tablerow} et blanc, aucun filet vertical.

\begin{center}
\small
\rowcolors{2}{tablerow}{white}
\begin{tabular}{@{}p{4.3cm}p{5.1cm}p{5.4cm}@{}}
\rowcolor{tableheader}
\textcolor{financeblue}{\textbf{Fichier}} &
\textcolor{financeblue}{\textbf{Contenu}} &
\textcolor{financeblue}{\textbf{Interdit}} \\
\code{<module de calcul>} & Logique métier, contrôles des données, calculs &
Toute primitive d'interface \\
\code{<point d'entrée>} & Saisie, déclenchement, agencement de l'interface &
Tout calcul métier \\
\code{<module d'affichage>} & Formatage des nombres, tableaux et tracés &
Tout calcul \\
\end{tabular}
\end{center}

\subsection{Tableau long en paysage}

Les tableaux longs utilisent \code{longtable} et les filets \code{booktabs},
en page paysage. Le tableau de la page suivante en montre la structure~:
en-tête répété, mention de suite, intertitres de groupe en gras bleu.

\begin{landscape}
\begin{center}
\scriptsize
\setlength{\tabcolsep}{4.5pt}
\renewcommand{\arraystretch}{1.22}
\begin{longtable}{@{}p{4.4cm}p{0.9cm}p{5.7cm}p{7.8cm}p{2.1cm}@{}}
\toprule
\textbf{Élément (fiche)} & \textbf{Réf.} & \textbf{Fonction(s) de calcul} &
\textbf{Fonction(s) / méthode de restitution} & \textbf{Clé} \\
\midrule
\endfirsthead
\toprule
\textbf{Élément (fiche)} & \textbf{Réf.} & \textbf{Fonction(s) de calcul} &
\textbf{Fonction(s) / méthode de restitution} & \textbf{Clé} \\
\midrule
\endhead
\multicolumn{5}{@{}l@{}}{\textit{Suite page suivante}} \\
\endfoot
\bottomrule
\endlastfoot

\multicolumn{5}{@{}l@{}}{\textbf{\textcolor{financeblue}{Groupe A~: intitulé du premier groupe}}} \\
\midrule
Premier élément & \ref{fiche:1} & \code{calcul\_a()} &
\code{restitution\_a()} (formule exacte) & \code{A1} \\
Deuxième élément & \ref{fiche:1} & \code{calcul\_b()} &
\code{restitution\_b()} (approximation, loi $\chi^2$) & aucune \\
\midrule
\multicolumn{5}{@{}l@{}}{\textbf{\textcolor{financeblue}{Groupe B~: intitulé du second groupe}}} \\
\midrule
Troisième élément & \ref{fiche:1} & \code{calcul\_c()} &
\code{restitution\_c()} (simulation) & \code{C1} \\
\end{longtable}
\end{center}
\end{landscape}


\partie{Partie III}{Fiches, figures et notes de mise en \oe uvre}

% =============================================================================
\section{Fiche méthode et fiche modèle}
\label{sec:fiche}
% =============================================================================

Chaque méthode du corps du document fait l'objet d'une fiche méthode en six
rubriques fixes, ouverte par l'environnement \code{fiche} et la macro
\code{\textbackslash Principe}\fsp; les rubriques suivantes sont appelées par
\code{\textbackslash Donnees}, \code{\textbackslash Conditions},
\code{\textbackslash Forces}, \code{\textbackslash Limites} et
\code{\textbackslash Referentiels}.

\begin{fiche}{Intitulé de la méthode}
\Principe
Ce que la méthode calcule et le modèle de référence qui la fonde, avec le
renvoi à sa fiche modèle (fiche~\ref{fiche:2}).
\Donnees
Données d'entrée, niveau de détail et profondeur d'historique nécessaires.
\Conditions
Hypothèses à vérifier sur les données avant l'emploi, et le moyen de les
vérifier.
\Forces
Ce qui justifie le choix de la méthode.
\Limites
Situations où la méthode se trompe, sens et ordre de grandeur de l'erreur.
\Referentiels
Référentiels où la méthode s'emploie (comptes sociaux, Solvabilité~II,
IFRS~17) et adaptations propres à chacun.
\end{fiche}

Chaque modèle cité fait l'objet, en annexe technique, d'une fiche modèle en
cinq rubriques fixes, référence en tête, ouverte par \code{\textbackslash
Cadre}\fsp; les rubriques suivantes sont appelées par \code{\textbackslash
Hypotheses}, \code{\textbackslash Utilisation}, \code{\textbackslash Formules}
et \code{\textbackslash LimitesModele}.

\begin{fiche}{Intitulé du modèle}
% Dans un document, ce \refl porte \src{ID} d'une référence de place du
% registre (contrôle 5) ; le registre n'en compte encore aucune de vérifiée.
\refl{A.~Auteur, \emph{Titre de la publication originale}, éditeur, année,
p.~10--12.}
\Cadre
Notations et objets du modèle~: un échantillon $x_1,\dots,x_n$.
\Hypotheses
Hypothèses du modèle, chacune avec l'effet de son non-respect.
\Utilisation
Méthodes du document qui emploient le modèle (fiche~\ref{fiche:1}) et
conditions de cet emploi.
\Formules
\[
  \hat\theta = \frac{1}{n}\sum_{i=1}^{n} x_i,
\]
avec, pour chaque formule, son statut épistémique~: résultat exact, résultat
asymptotique, approximation ou constat empirique.
\LimitesModele
Ce que le modèle ne représente pas, et renvoi aux pièges de la fiche méthode.
\end{fiche}

% =============================================================================
\section{Figure légendée}
\label{sec:figure}
% =============================================================================

Les figures sont placées à l'endroit exact de l'appel par l'environnement
\code{figureH} (et \code{tableH} pour un tableau légendé)\fsp; les environnements
flottants \code{figure} et \code{table} restent disponibles. L'étiquette de
légende est en gras, en petites capitales et en couleur \code{maincolor}\fsp;
une légende d'une ligne est centrée, une légende plus longue est justifiée
(figure~\ref{fig:exemple}).

\begin{figureH}
\centering
\fcolorbox{financeblue}{tablerow}{\begin{minipage}[c][4cm][c]{9cm}
\centering\small\color{financeblue} Emplacement de la figure\\
\code{\textbackslash includegraphics[width=...]\{fichier.pdf\}}
\end{minipage}}
\caption{Exemple de légende longue. Le texte de la légende est composé en
romain au corps du texte, justifié sur toute la largeur\fsp; seule l'étiquette
porte la couleur de la charte.}
\label{fig:exemple}
\end{figureH}

\begin{figureH}
\centering
\fcolorbox{financeblue}{tablerow}{\begin{minipage}[c][2cm][c]{9cm}
\centering\small\color{financeblue} Emplacement de la figure
\end{minipage}}
\caption{Légende courte, centrée.}
\end{figureH}

% =============================================================================
\section{Notes de mise en \oe uvre}
\label{sec:notes}
% =============================================================================

\subsection{Packages remplacés}

\begin{center}
\small
\rowcolors{2}{tablerow}{white}
\begin{tabular}{@{}p{3.0cm}p{11.8cm}@{}}
\rowcolor{tableheader}
\textcolor{financeblue}{\textbf{Package retiré}} &
\textcolor{financeblue}{\textbf{Remplacement}} \\
\code{newtxtext} & \code{tgtermes} (repli sur \code{ptm}) pour le texte,
\code{txtt} pour le code \\
\code{babel-french} & Césure française, \code{\textbackslash frenchspacing},
libellés, date, retrait du premier paragraphe, puces et espacement des
listes, réglés dans le préambule \\
\code{csquotes} & Guillemets «~» saisis directement (codage T1), avec espaces insécables \\
\code{caption} & Redéfinition de \code{\textbackslash @makecaption} \\
\code{float} & Environnements \code{figureH} et \code{tableH} (bloc non flottant) \\
\code{enumitem} & \code{itemize} et \code{enumerate} redéfinis, clés
\code{itemsep}, \code{topsep}, \code{parsep}, \code{leftmargin}\fsp; liste
\code{rubriques} pour les fiches \\
\code{mdframed} & \code{framed}, avec cadre sur mesure \\
\code{titlesec} & \code{\textbackslash @startsection} et
\code{\textbackslash @seccntformat} \\
\code{tocloft} & Redéfinition de \code{\textbackslash l@section} \\
\code{pdflscape} & \code{lscape} et \code{/Rotate 90} ajouté aux pages paysage (\code{atbegshi}) \\
\code{tikz} & Filets de la page de garde superposés à la page par
\code{atbegshi} (sans crochet du noyau récent) \\
\code{dsfont}, \code{mathrsfs} & Non utilisés\fsp; \code{\textbackslash mathbb}
(amssymb) reste disponible \\
\end{tabular}
\end{center}

\subsection{Ponctuation}

Sans babel-french, l'espace qui précède les signes doubles n'est plus
ajoutée automatiquement, et une espace ordinaire saisie avant eux autorise
une coupure de ligne malvenue. On saisit donc \code{\textasciitilde:} pour le
deux-points (espace insécable) et \code{\textbackslash fsp;},
\code{\textbackslash fsp!}, \code{\textbackslash fsp?} pour les autres
signes (espace fine insécable)\fsp; les guillemets se saisissent
directement, avec une espace insécable à l'intérieur~:
\code{«\textasciitilde{}texte\textasciitilde{}»} donne «~texte~».

\end{document}
